\section{Second-Order Training Objective}
\label{app:second-order-training}

The second equality in \eqref{eq:second-loss} follows from
$F_{\theta,\mathbb R}-\mathbf h_{\mathbb R}
=\sqrt{v_{\mathbb R}}\mathbf b$ and $2Dv_{\mathbb R}=Dv$.

For the fixed mean estimator $F_\theta$, minimizing over scalar
functions gives the population target
\begin{equation}
 \tau_\theta(R,v)=\frac1D
 \E[\norm{F_\theta(R,v)-H}_F^2\mid R,v,\alpha=1].
 \label{eq:conditional-risk}
\end{equation}
Indeed, with $U=\norm{F_\theta-H}_F^2/(Dv)$, conditioning
\eqref{eq:second-loss} on $(R,v,\alpha=1)$ gives
$(a_\psi-\E[U\mid R,v,\alpha=1])^2+
\operatorname{Var}(U\mid R,v,\alpha=1)$.
Moreover, expanding the squared error about the conditional mean gives
\begin{equation}
 \tau_\theta(R,v)=\tau_{\rm C}(R,v)
 +\frac{\norm{F_\theta(R,v)-m(R,v)}_F^2}{D},
 \label{eq:risk-decomposition}
\end{equation}
where $\tau_{\rm C}=D^{-1}
\E[\norm{H-m}_F^2\mid R,v,\alpha=1]$.
The scalar output in \eqref{eq:second-readout} therefore learns the
conditional mean-squared error of the deployed mean estimator. This
target reduces to the posterior covariance trace per complex coordinate
when $F_\theta=m$.

To derive the average in \eqref{eq:pooled-variance}, fit Gaussian
distributions with fixed means $F_k$ and covariance $\tau\I_D$ for
every terminal. Giving equal weight to each terminal's channel
distribution conditioned on being active yields
the expected negative log density, up to a constant,
\begin{equation}
 \mathcal J(\tau)=KD\log\tau+
 \frac{D}{\tau}\sum_{k=1}^{K}\tau_\theta(R_k,v).
 \label{eq:shared-error-objective}
\end{equation}
Its derivative is
$KD(\tau-K^{-1}\sum_k\tau_\theta(R_k,v))/\tau^2$.
The minimizer under $\tau\geq\tau_{\min}$ is the average conditional
error clipped below by $\tau_{\min}$.
Equation~\eqref{eq:pooled-variance} approximates this fit with the learned
error estimates. Equal weighting concerns each terminal's distribution
conditioned on being active; activity probabilities enter later through
the mixture in \eqref{eq:mixture}.
